% PROGRAMA DE PÓS-GRADUAÇÃO EM ECONOMIA APLICADA
% UNIVERSIDADE FEDERAL DO RIO GRANDE - FURG

% ===========================================================

\documentclass[a4paper, 17pt]{extarticle}

%%%%%%%%%%%%%%%%%%%%%%%%% PACOTES %%%%%%%%%%%%%%%%%%%%%%%%%%%%%%%%%%
 
\usepackage[margin = 0.5in, bottom = 2in]{geometry}
\usepackage[utf8]{inputenc}
\usepackage[T1]{fontenc}
\usepackage{graphicx}
\usepackage{enumitem}
\usepackage{multicol}
\usepackage{subcaption}
\usepackage{arydshln}
%\usepackage{pdfcomment}
\usepackage{parskip}
\usepackage{xcolor}
\usepackage{lipsum}	
\usepackage{amsfonts,amsmath,amssymb, amsthm, pifont}
\usepackage{booktabs}
\usepackage[colorlinks=true, linkcolor=blue, citecolor=blue, urlcolor=blue]{hyperref} 
\usepackage{todonotes} 
\usepackage{dsfont}
\usepackage{cancel}
\usepackage{layout}
\usepackage[most]{tcolorbox}
\usepackage{tikz}
\usepackage{tikz-3dplot} 
\usepackage[absolute, overlay]{textpos} 
\usetikzlibrary{arrows.meta, cd, patterns, patterns.meta}
\usepackage[bottom]{footmisc}

% Configuração da Fonte
\usepackage{times} % Fonte: Times New Roman

% PACOTES ESSENCIAIS
\usepackage{caption}  
\usepackage{titlesec}  

\titleformat{\section}
  {\normalfont\Large\bfseries} % formato de la fuente
  {\S \hspace{-0cm} \thesection}              % etiqueta de la sección (con símbolo ♦)
  {0.5em}                        % espacio entre el número y el título
  {}                           % código antes del título

%\setlength{\footskip}{10pt} % Espacio entre texto y pie de página
\renewcommand{\footnoterule}{\vspace{5pt}\hrule width 0.3\linewidth\vspace{5pt}}
\setlength{\textheight}{750pt}
%\setlength{\textwidth}{530pt}
%\setlength{\headsep}{10pt}
%\setlength{\marginparwidth}{145pt}
%\setlength{\marginparsep}{12pt}
%\setlength{\topmargin}{-60pt}

\definecolor{verdeoscuro}{RGB}{0,100,0} % Un verde bastante oscuro
\definecolor{rojoscuro}{RGB}{139,0,0} % rojo oscuro tipo "DarkRed"

%\pagestyle{empty}

\newcommand{\ejemplo}[1]{%
  \begin{tcolorbox}[colframe=verdeoscuro, colback=white!1, boxrule=0.8pt]
    #1
  \end{tcolorbox}
}
\newcommand{\alerta}[1]{%
  \begin{tcolorbox}[colframe=red, colback=white!1, boxrule=0.8pt]
    #1
  \end{tcolorbox}
}
\newcommand{\demo}[1]{%
  \begin{tcolorbox}[colframe=gray, colback=white!1, boxrule=0.8pt]
    \textcolor{gray}{#1}
  \end{tcolorbox}
}

\todostyle{orange}{
    linecolor=orange,        % Flecha negra
    bordercolor=orange,      % Borde negro
    textcolor=black!90,      % Borde negro
    backgroundcolor=white!33,  % Sin relleno de color
    size=\scriptsize           % Tamaño de texto
}
\todostyle{green}{
    linecolor=verdeoscuro,        % Flecha negra
    bordercolor=verdeoscuro,      % Borde negro
    textcolor=black!90,      % Borde negro
    backgroundcolor=white!33,  % Sin relleno de color
    size=\scriptsize           % Tamaño de texto
}
\todostyle{red}{
    linecolor=red,        % Flecha negra
    bordercolor=red,      % Borde negro
    textcolor=black!90,      % Borde negro
    backgroundcolor=white!33,  % Sin relleno de color
    size=\scriptsize           % Tamaño de texto
}
\todostyle{red2}{
  linecolor=rojoscuro,        % Flecha negra
  bordercolor=rojoscuro,      % Borde negro
  textcolor=black!90,      % Borde negro
  backgroundcolor=white!33,  % Sin relleno de color
  size=\scriptsize           % Tamaño de texto
}
\todostyle{gris}{
  linecolor=gray,        % Flecha negra
  bordercolor=gray,      % Borde negro
  textcolor=black!90,      % Borde negro
  backgroundcolor=white!33,  % Sin relleno de color
  size=\scriptsize           % Tamaño de texto
}
\begin{document} %Início do Documento

%\layout
%\newpage 


%%%%%%%%%%%%%%%%%%%%%%%%%  MEU PREAMBULO %%%%%%%%%%%%%%%%%%%%%%%%%%%%%%%%%

\makeatletter
\renewcommand\@makefnmark{%
  \hbox{\@textsuperscript{\color{magenta}\normalfont\@thefnmark}}}
\renewcommand\@makefntext[1]{%
  \noindent\makebox[1.8em][r]{\textsuperscript{\color{magenta}\@thefnmark}}#1}
\makeatother

\theoremstyle{definition}
\newtheorem*{definition}{\rotatebox{45}{\large \(\square\)}}

\theoremstyle{plain}
\newtheorem*{lemma}{{\scriptsize \(\square\)}}
\newtheorem*{proposition}{{\large \(\square\)}}
\newtheorem*{theorem}{{\large \(\blacksquare\)}} 

\makeatletter
% Eliminar la puntuación del estilo 'plain' y 'definition'
\def\@thm@headpunct{} % Elimina el punto de los títulos de los teoremas

% Modificar el estilo de 'plain' y 'definition'
\def\th@plain{%
  %\thm@headfont{\normalfont\bfseries} % Mantener la fuente normal, sin negrita
  %\thm@notefont{\normalfont} % El texto de la nota (si la hay) no será en negrita
  \thm@headpunct{} % Sin puntuación adicional 
}

\def\th@definition{%
  %\thm@headfont{\normalfont\bfseries} % Fuente normal para el título
  %\thm@notefont{\normalfont} % El texto de la nota (si la hay) no será en negrita
  \thm@headpunct{} % Sin puntuación adicional
}

\def\th@remark{%
  %\thm@headfont{\normalfont\bfseries} % Fuente normal para el título
  %\thm@notefont{\normalfont} % El texto de la nota (si la hay) no será en negrita
  \thm@headpunct{} % Sin puntuación adicional
}
\makeatother



\theoremstyle{remark}
\newtheorem*{example}{\textcolor{verdeoscuro}{\underline{{Exemplo}}}}
\newtheorem*{exercise}{\textcolor{rojoscuro}{\underline{{Exercise}}}}\theoremstyle{remark}
\newtheorem*{note}{\textnormal{\emph{Nota. }}}
\newtheorem*{corollary}{\textnormal{\emph{Corolário. }}}

\newcommand{\id}{\mathds{1}}
\newcommand{\R}{\mathbb{R}}
\newcommand{\B}{\mathbb{B}}
\newcommand{\Ei}{}
\newcommand{\Ef}{}
\newcommand{\N}{\mathbb{N}}
\newcommand{\U}{\mathcal{U}}
\newcommand{\V}{\mathcal{V}}
\newcommand{\E}{\mathbb{S}}
\newcommand{\ab}{^{\text{\ding{71}}}}
\newcommand{\ce}{^{\text{\ding{70}}}}
\newcommand{\x}{\textbf{x}}
\newcommand{\y}{\textbf{y}}


\setlength{\columnsep}{1in}

%\begin{multicols}{2}

\-\vspace{-1.3cm} 

\par\noindent \includegraphics[width = \linewidth]{Imagens/Header.png} 

%%%%%%%%%%%%%%%%%%%%%%%% CONFIG. DO TÍTULO %%%%%%%%%%%%%%%%%%%%%%%%

\begin{center} %Inicia a centralização
{%{ \Large QUALI SURVIVAL GUIDE COLLECTION} \\ 
\LARGE \textbf{\(^{+3}\)Notas de Topologia}} \vspace{0.1in}   \\
 { \today }\\ 
 {\small
 \vspace{0.4cm}\emph{"Alguma frase" }\hfill Nestor Heli Aponte Avila\(^{1}\) \\ 
  
Dono da frase \hfill \href{mailto:n267452@dac.unicamp.br}{\url{n267452@dac.unicamp.br}}} %\\
%Prof. Piper Pimienta \\
%\href{mailto:pipe@dac.unicamp.br}{\url{pipe@dac.unicamp.br}}}
\end{center} %Finaliza a centralização
\- \vspace{-1cm}

{\small ** Abstrasct muito abstract . **\vspace{0.3cm}}%, destinado àqueles que já possuem alguma familiaridade com o assunto, pois não contém demonstrações formais, apenas as ideias por trás delas. Sinceramente, um fã da simplicidade. \\ 


%%%%%%%%%%%%%%%%%%%%%%%%% CONTENIDO %%%%%%%%%%%%%%%%%%%%%%%%%%%%%%%%%%
{\large \textbf{Notação}} \hfill \text{\rotatebox{45}{\large $\square$}}  Definição \hfill \((-)^{\text{\ding{71}}}\)  Aberto \hfill  \((-)^{\text{\ding{70}}}\) Fechado \\ 
\- \hspace{4.7cm}{\scriptsize$\square$} \ Lema \hspace{2.9cm} {\large$\square$} \ Proposição \hspace{1.3cm}  {\large $\blacksquare$}  \ Teorema

\section{Conjuntos}

\begin{lemma}[Leis de De Morgan]
    Se \(A_i \subset X\) então, 
    \[(a)\ X\setminus \bigcup A_i = \bigcap (X \setminus A_i) \ \ \text{ e } \ \ (b) \ X\setminus \bigcap A_i = \bigcup (X\setminus A_i). \]
    \textcolor{gray}{\(\rightarrow\) Faz \((a)\) por double contenção e use ele para concluir \((b)\)}
\end{lemma}

\begin{lemma}[Pre-imagen bém portada]
    Sejam \(f:X\to Y\) função e \(B_i \subset Y\). Então,
    \[(a)\ f^{-1}\left(\bigcup B_i\right) = \bigcup f^{-1}(B_i); \ \ \ (b)\ f^{-1}\left(\bigcap B_i\right) = \bigcap f^{-1}(B_i)\]
    \[(c)\ f^{-1}\left(Y \setminus B_i\right) = X \setminus f^{-1}(B_i)\]
\end{lemma}


  

%%%%%%%%%%%%%%%%%%%%%%%%%% REFERÊNCIAS %%%%%%%%%%%%%%%%%%%%%%%%%%%%%%%%

\begin{thebibliography}{9}

%\bibitem{martin1966complex}
%Martin, D. y Ahlfors, L.V. (1966). \textit{Complex Analysis}. New York: McGraw-Hill.

\bibitem{lima2004analise}
Lima, Elon Lages (2004). \textit{An{\'a}lise real Vol. 2}. IMPA, Rio de Janeiro.

\end{thebibliography}

%\end{multicols}

%%%%%%%%%%%%%%%%%%%%%%%%%% APÊNDICES %%%%%%%%%%%%%%%%%%%%%%%%%%%%%%%%

\end{document}
